\documentclass[12pt]{article}

\usepackage[utf8]{inputenc}
\usepackage[T2A]{fontenc}
\usepackage[ukrainian]{babel}
\usepackage{geometry}
\geometry{a4paper, margin=2cm}

\usepackage{xcolor}
\usepackage{tcolorbox}
\usepackage{graphicx}
\usepackage{amsmath, amsthm, amssymb}

\usepackage{caption}
\captionsetup{labelsep=period}

\newtheorem{theorem}{Теорема}

\usepackage{fancyhdr}
\pagestyle{fancy}
\fancyhf{}
\setlength{\headheight}{15pt}

\newcommand{\student}{Левченко Анна}
\newcommand{\labdate}{\today}

\fancyhead[L]{\student}
\fancyhead[R]{\labdate}
\fancyfoot[C]{\thepage}

\begin{document}

\begin{center}
    \Large \textbf{Лабораторна робота № 9} \\
    \large \textbf{Тема: Набір математичних текстів та автоматизація в LaTeX}
\end{center}

\section*{Мета роботи}
\textit{навчитися вводити математичні формули та структурувати документ,
використовувати кольорові рамки для визначень, теорем та прикладів,
а також вставляти зображення}.

\section{Короткі теоретичні відомості}

\begin{tcolorbox}[colback=blue!10!white, colframe=red!80!white, title=Визначення]
Квадратом називають прямокутник, у якого всі сторони рівні.
\end{tcolorbox}

На рисунку~\ref{fig:square} зображено квадрат $ABCD$. Прямокутник є
паралелограмом, тому і квадрат~--- паралелограм, у якого всі сторони рівні,
тобто він є і ромбом. Отже, квадрат має властивості прямокутника та ромба.

\begin{figure}[h]
    \centering
    \includegraphics[width=0.3\textwidth]{foto1.png}
    \caption{Квадрат $ABCD$}
    \label{fig:square}
\end{figure}

\medskip
\textit{Властивості квадрата:}
\begin{enumerate}
    \item Усі кути квадрата прямі.
    \item Периметр $P_{ABCD} = 4AB$.
    \item Діагоналі квадрата між собою рівні (рис.~\ref{fig:square}).
    \item Діагоналі квадрата взаємно перпендикулярні й точкою перетину
          діляться навпіл (рис.~\ref{fig:diag}).
    \item Діагоналі квадрата ділять його кути навпіл, тобто утворюють кути
          $45^\circ$ зі сторонами квадрата (рис.~\ref{fig:diag}).
    \item Точка перетину діагоналей квадрата рівновіддалена від усіх його
          вершин: $AO = BO = CO = DO$.
\end{enumerate}

\begin{figure}[h]
    \centering
    \includegraphics[width=0.3\textwidth]{foto2.png}
    \caption{Діагоналі квадрата}
    \label{fig:diag}
\end{figure}
\newpage

\begin{tcolorbox}[colback=green!5!white, colframe=green!75!black, title=Приклад 1]
Точки $K$ і $M$ належать відповідно діагоналям $BD$ і $AC$ квадрата $ABCD$,
причому $AM = \dfrac{1}{4}AC$, $BK = \dfrac{1}{4}BD$.
Доведіть, що $\triangle ADM = \triangle BAK$ (рис.~\ref{fig:example}).
\end{tcolorbox}

\begin{figure}[h]
    \centering
    \includegraphics[width=0.3\textwidth]{foto3.png}
    \caption{Ілюстрація до прикладу}
    \label{fig:example}
\end{figure}

\begin{proof}
\begin{enumerate}
    \item $\angle MAD = \angle ABK = 45^\circ$, а $AD = AB$ як сторони квадрата.
    \item Оскільки $AC = BD$ як діагоналі квадрата і
          $AM = \dfrac{1}{4}AC$, $BK = \dfrac{1}{4}BD$, то $AM = BK$.
    \item Отже, $\triangle AMD = \triangle BKA$ за двома сторонами та кутом між ними.
\end{enumerate}
\end{proof}

\section{Ознаки квадрата}

\begin{tcolorbox}[colback=blue!5!white, colframe=blue!75!black, title=Ознаки квадрата]
\begin{theorem}
\textit{1) Якщо діагоналі прямокутника взаємно перпендикулярні, то він є
квадратом.} \\
\textit{2) Якщо діагоналі ромба між собою рівні, то він є квадратом.}
\end{theorem}
\end{tcolorbox}

\begin{proof}
1) Прямокутник є паралелограмом, а паралелограм із взаємно перпендикулярними
діагоналями є ромбом. Отже, у заданого прямокутника всі сторони рівні,
а тому він є квадратом.

2) Ромб є паралелограмом, а паралелограм з рівними діагоналями є
прямокутником. Отже, у розглянутому ромбі всі кути прямі, а тому він є
квадратом.
\end{proof}

\section*{Контрольні запитання}
\begin{enumerate}
    \item Що таке квадрат? Які його властивості?
    \item Які ознаки квадрата через діагоналі?
    \item Як вставляти зображення та підпис у LaTeX?
    \item Як оформити визначення, теорему та приклад у кольоровій рамці?
\end{enumerate}

\end{document}
